% META: {"id": "TSPE-LIMFCT-ME-002", "chapitre": "TSPE-LIMITES-FONCTIONS", "type_objet": "methode", "capacites_codes": ["C2"], "methodes": ["M2"], "sources_inspiration": [], "status": "approved"}
\begin{fichemethode}{M2}{Déterminer les asymptotes d'une courbe}

\quandUtiliser{%
  Utiliser cette méthode lorsque l'énoncé demande de « déterminer les asymptotes », « interpréter graphiquement les limites » ou « étudier le comportement aux bornes du domaine ».
}

\pasApas{%
  \begin{enumerate}
    \item \textbf{Déterminer le domaine de définition} et identifier les valeurs interdites.
    \item \textbf{Asymptotes horizontales :} calculer $\lim_{x \to +\infty} f(x)$ et $\lim_{x \to -\infty} f(x)$. Si une limite vaut $\ell$ fini, alors $y = \ell$ est AH.
    \item \textbf{Asymptotes verticales :} pour chaque valeur interdite $a$, calculer $\lim_{x \to a^+} f(x)$ et $\lim_{x \to a^-} f(x)$. Si une limite est $\pm\infty$, alors $x = a$ est AV.
    \item \textbf{Preciser la position de la courbe :} si besoin, indiquer si $f(x) > \ell$ ou $f(x) < \ell$ au voisinage de l'AH, et de quel cote la courbe se situe par rapport a l'AV.
  \end{enumerate}
}

\exempleRedige{%
  \textbf{Exemple.} Déterminer les asymptotes de $f(x) = \dfrac{3x - 1}{x + 2}$.

  \medskip
  \commentaireMarge{Domaine : $\mathbb{R} \setminus \{-2\}$.}%
  \textbf{Domaine :} $f$ est définie pour $x
\neq -2$.

  \commentaireMarge{AH : factoriser par $x$.}%
  \textbf{AH :} $f(x) = \dfrac{x(3 - 1/x)}{x(1 + 2/x)} = \dfrac{3 - 1/x}{1 + 2/x} \to \dfrac{3}{1} = 3$ en $\pm\infty$.

  La droite $y = 3$ est AH en $+\infty$ et en $-\infty$.

  \commentaireMarge{AV : étudier les limites en $-2$.}%
  \textbf{AV :} Numérateur en $-2$ : $3(-2) - 1 = -7
\neq 0$. Dénominateur $\to 0$.
  \begin{itemize}
    \item $x \to (-2)^+$ : $x + 2 \to 0^+$, $f(x) \to \dfrac{-7}{0^+} = -\infty$.
    \item $x \to (-2)^-$ : $x + 2 \to 0^-$, $f(x) \to \dfrac{-7}{0^-} = +\infty$.
  \end{itemize}

  La droite $x = -2$ est AV.
}

\pieges{%
  \begin{itemize}
    \item \textbf{Piège 1.} Chercher une AV en un point ou $f$ est définie et continue.
    \item \textbf{Piège 2.} Oublier d'étudier separement les limites a gauche et a droite.
    \item \textbf{Piège 3.} Affirmer sans justification que la courbe « ne coupe pas » l'asymptote.
  \end{itemize}
}

\verifier{%
  Sur la calculatrice graphique, tracer la fonction et vérifier visuellement la presence des asymptotes trouvees.
}

\sEntrainer{\refExos{M2}}

\end{fichemethode}
